\scanpage{58}
Soit $\Pi_n$ ($n\in\mathbb N$, $2\leq n\leq+\infty$) le polygône
combinatoire défini par $\unsure{D_n|\mathbb D_n}$, considéré comme
espace homogène sous $G_1$. Alors
\[
\Pi_n=(S_n,A_n,R_n)
\]
avec
\[
\underset{\substack{\Vert\\\{s_i\}_{i\in\mathbb Z_n}}}{S_n}=\deleted{\unsure{D_n/\sigma_1|D_n}}\underset{(\overline{\sigma_0})}{u\sigma}\backslash D_n\simeq\mathbb Z_n=
\]
[car $(u\sigma{\unsure{]|)}}^2=1$ \unsure{ou|or|donc} $u\sigma$ définit une section
de l'ext. de $\mathbb Z_2$ par $\mathbb Z_n$]
\[
\underset{\substack{\Vert\\\{a_i\}_{i\in\mathbb Z_n}}}{A_n}=\underset{\substack{\Vert\\(\sigma_1)}}{\sigma}\backslash D_n\simeq\mathbb Z_n
\]
Enfin, la relation d'incidence est donnée
par \qquad $(s_i,a_i)$ et $(s_{i+1},a_i)$ incidents.
\figure{\begin{tikzpicture}[x=1cm,y=1cm,baseline]
\draw (0,0) -- (1.55,.09) -- (2.95,1.04) -- (3.03,2.02);
\fill (1.55,.09) circle (.04) (2.95,1.04) circle (.04);
\node[below left,inner sep=1pt] at (0,0) {$s_0$};
\node[above,inner sep=1pt] at (.69,.09) {$a_0$};
\node[below right,inner sep=1pt] at (1.55,.09) {$s_1$};
\node[above left,inner sep=2pt] at (2.26,.65) {$a_1$};
\node[right,inner sep=2pt] at (2.95,1.04) {$s_2$};
\end{tikzpicture}}
On retrouve la description des polyg. comb.
\unsure{standart|standard|d'avant}, et le fait que la classe d'iso est
caractérisée par le nb des côtés = nb d'arêtes

\uline{\unsure{(Re)corollaire|Corollaire} 2} Si $\Pi,\Pi'$ sont deux polygônes
combinatoires iso, \deleted{alors} et $r\in\mathrm{Rep}(\Pi)$,
alors
\diagram{\begin{tikzcd}[ampersand replacement=\&,row sep=small,column sep=large]
\mathrm{Isom}(\Pi,\Pi')\arrow[r,"\sim"]\&\mathrm{Rep}(\Pi')\\
\varphi\arrow[r,mapsto]\&\varphi(r)
\end{tikzcd}}
C'est un cas particulier d'un \unsure{énoncé|théorème} sur les
% archive sheet number: 57
